Let $p - q \equiv 1 \pmod 4$, $A = \Cl(p,q)$, and let $\nu$ be one adjoined generator orthogonal to the generating space with $\nu^2 = \gamma \in \{\pm 1\}$, so that $B_\gamma = \Cl(p{+}1,q)$ or $\Cl(p,q{+}1)$. Then: conjugation by $\nu$ restricts on $A$ to the grade involution, $\nu a \nu^{-1} = \hat{a}$; the map $a \mapsto \varepsilon_+ a$ is an isomorphism $\Cl^0(p,q) \to \varepsilon_+ A$; the elements $E_{11} = \varepsilon_+$, $E_{22} = \varepsilon_-$, $E_{12} = \varepsilon_+\nu$, $E_{21} = \gamma\varepsilon_-\nu$ form a full $2 \times 2$ system of matrix units in $B_\gamma$; hence $B_\gamma \cong \Mat_2(\Cl^0(p,q))$ is simple, the images of $\varepsilon_\pm$ are a conjugate pair of non-central idempotents exchanged by $\Ad_\nu$, and the two ascents are isomorphic as algebras though non-isometric as quadratic data: the sign $\gamma$ is a geometric datum with no algebraic shadow. One nondegenerate generator is minimal, the null generator never fuses, and fusion fires at most once in any transport history. \emph{Provenance:} \cite[Prop.~5.2, Lemma 5.3, Lemma 5.4, Thm.~5.5, Thm.~5.7, Remark 5.8, Prop.~7.4]{ButlerT1}.
